\chapter{Diseases of the Bones and Joints}
\label{ch:bones}
The human skeleton consists of 206 bones connected by joints, providing structural support, movement, mineral storage, and hematopoiesis. Joint diseases are among the most common causes of disability worldwide. This chapter covers metabolic bone diseases, inflammatory arthropathies, degenerative joint diseases, infectious conditions, and bone neoplasms.

%-------------------------------------------------------------
\section{Bone Infections}
%-------------------------------------------------------------

%-------------------------------------------------------------

%-------------------------------------------------------------

%-------------------------------------------------------------

Bone infections, predominantly osteomyelitis and septic arthritis, represent severe inflammatory conditions resulting from microbial invasion of osseous architecture or synovial spaces. Infection occurs via hematogenous seeding, direct inoculation from trauma or surgery, or contiguous spread from adjacent soft tissue infections. Inadequate or delayed treatment can lead to cortical bone necrosis, chronic sinus tract formation, and permanent structural or articular disability.

%-------------------------------------------------------------

%-------------------------------------------------------------

%-------------------------------------------------------------

%-------------------------------------------------------------

\label{sec:Bone Infections}
%-------------------------------------------------------------%-------------------------------------------------------------

\subsection{Osteomyelitis}
\label{sec:Osteomyelitis}
Osteomyelitis is infection of bone, which may be acute or chronic. Hematogenous osteomyelitis (most common in children) typically affects the metaphysis of long bones (femur, tibia, humerus), with \textit{Staphylococcus aureus} as the predominant organism. Contiguous-focus osteomyelitis occurs after direct inoculation (trauma, surgery, open fractures) or spread from adjacent soft tissue infection. The condition results from bacterial infection of bone leading to inflammation, necrosis, and new bone formation. Clinical features in children include localized bone pain, fever, and elevated inflammatory markers; in adults, it more commonly involves the vertebrae (vertebral osteomyelitis/diskitis). Chronic osteomyelitis is characterized by bone necrosis (sequestrum), new bone formation (involucrum), and sinus tracts draining pus. Diagnosis involves blood cultures, inflammatory markers (ESR, CRP), and imaging (MRI is the gold standard for detecting bone edema and abscesses; CT for sequestrum; bone scan). Management requires prolonged antibiotic therapy (6 weeks or more, IV initially then oral step-down) and surgical debridement for chronic or refractory disease. Prognosis is favorable with early diagnosis and appropriate treatment.

\subsection{Septic Arthritis}
\label{sec:Septic Arthritis}
Septic arthritis is infection of a joint, a medical emergency requiring prompt treatment to prevent irreversible articular cartilage destruction. It most commonly affects large joints (knee, hip) and is usually monoarticular. Routes of infection include hematogenous spread (most common), direct inoculation (trauma, surgery, injection), and contiguous spread from adjacent osteomyelitis. \textit{S.~aureus} is the most common organism in adults; \textit{Neisseria gonorrhoeae} in young, sexually active adults; and \textit{Kingella kingae} in children under 4. The condition results from bacterial infection of the joint space. Clinical features include acute joint pain, swelling, warmth, erythema, and reduced range of motion, with systemic signs of infection (fever, chills). Diagnosis requires joint aspiration showing turbid fluid with elevated WBC count (>50,000 cells/$\mu$L with >90\% neutrophils), positive Gram stain and culture. Management involves urgent joint drainage (repeated aspiration, arthroscopic lavage, or open drainage) and IV antibiotics (empiric: vancomycin + ceftriaxone, adjusted based on culture and sensitivity). Prognosis is favorable with prompt treatment; delayed treatment leads to irreversible joint damage.

%-------------------------------------------------------------
\section{Bone Tumors}
%-------------------------------------------------------------

%-------------------------------------------------------------

%-------------------------------------------------------------

%-------------------------------------------------------------

Bone tumors encompass a spectrum of benign and malignant primary neoplasms arising from osteogenic, chondrogenic, fibrogenic, or hematopoietic cell lines within bone architecture. Primary malignant bone tumors disproportionately affect children and young adults during periods of rapid skeletal growth, whereas metastatic bone lesions predominate in older individuals. Definitive diagnosis requires integration of radiographic, histological, and molecular features to guide limb-sparing surgery and systemic therapy.

%-------------------------------------------------------------

%-------------------------------------------------------------

%-------------------------------------------------------------

%-------------------------------------------------------------

\label{sec:Bone Tumors}
%-------------------------------------------------------------%-------------------------------------------------------------

\subsection{Ewing Sarcoma}
\label{sec:Ewing Sarcoma}
Ewing sarcoma is the second most common primary bone malignancy in children and adolescents, arising from primitive neuroectodermal cells. The classic translocation t(11;22)(q24;q12) results in the EWS-FLI1 fusion gene. It typically affects the diaphysis of long bones (femur, tibia, fibula) and the pelvis. The condition results from the EWS-FLI1 fusion gene driving oncogenesis. Clinical features include bone pain, swelling, and a mass, with systemic symptoms (fever, weight loss) that may mimic infection. X-rays show a permeative lytic lesion with "onion skin" periosteal reaction. MRI and CT chest are essential for staging. Diagnosis is confirmed by biopsy with molecular testing for the translocation. Management includes neoadjuvant chemotherapy (vincristine, doxorubicin, cyclophosphamide alternating with ifosfamide/etoposide), local control with surgery and/or radiation therapy, and adjuvant chemotherapy. Prognosis is favorable for localized disease; the 5-year survival rate is approximately 70\%.

\subsection{Multiple Myeloma (Bone Manifestations)}
\label{sec:Multiple Myeloma (Bone Manifestations)}
Multiple myeloma is discussed in detail in Chapter~\ref{ch:blood}. Bone disease is a hallmark feature, present in over 80\% of patients at diagnosis. Myeloma cells produce factors (RANKL, MIP-1$\alpha$) that activate osteoclasts, leading to lytic bone lesions without osteoblastic repair. The most common sites are the spine, skull, ribs, pelvis, and long bones. Clinical features include pathologic fractures, pain, vertebral compression fractures with height loss and kyphosis, and spinal cord compression. Diagnosis of bone disease involves skeletal survey (low sensitivity), CT (high sensitivity for lytic lesions), MRI (detects marrow infiltration), and PET-CT. Management includes bisphosphonates (zoledronic acid, pamidronate) or denosumab to reduce skeletal events, radiation for symptomatic lesions, kyphoplasty/vertebroplasty for compression fractures, and myeloma-directed therapy.

\subsection{Osteosarcoma}
\label{sec:Osteosarcoma}
Osteosarcoma is the most common primary malignant bone tumor (excluding myeloma), with a bimodal age distribution (peak in adolescents during the growth spurt and in older adults). It arises from mesenchymal cells producing osteoid and most commonly affects the metaphysis of long bones (distal femur, proximal tibia, proximal humerus). The condition results from malignant transformation of osteoblast precursors; risk factors include prior radiation, Paget's disease, Li-Fraumeni syndrome, hereditary retinoblastoma, and Bloom syndrome. Clinical features include localized bone pain, swelling, and a palpable mass. X-rays show a destructive lesion with sunburst periosteal reaction and Codman's triangle. Diagnosis is confirmed by biopsy. Management involves neoadjuvant chemotherapy (doxorubicin, cisplatin, high-dose methotrexate), surgical resection (limb-salvage surgery when possible), and adjuvant chemotherapy. Prognosis depends on tumor necrosis after neoadjuvant chemotherapy (>90\% necrosis is a favorable prognostic factor).

%-------------------------------------------------------------
\section{Degenerative Joint Disease}
%-------------------------------------------------------------

%-------------------------------------------------------------

%-------------------------------------------------------------

%-------------------------------------------------------------

%-------------------------------------------------------------

%-------------------------------------------------------------

%-------------------------------------------------------------

%-------------------------------------------------------------

\label{sec:Degenerative Joint Disease}
%-------------------------------------------------------------%-------------------------------------------------------------

\subsection{Osteoarthritis}
\label{sec:Osteoarthritis}
Osteoarthritis (OA) is the most common joint disease worldwide, characterized by progressive articular cartilage degradation, subchondral bone remodeling, osteophyte formation, and synovial inflammation. It affects over 500 million people globally, with prevalence increasing with age. The condition results from mechanical and metabolic factors leading to cartilage breakdown, with risk factors including advancing age, obesity, joint injury, female sex, and genetic predisposition. Clinical features include joint pain (worse with activity and relieved by rest), stiffness (morning stiffness <30 minutes, unlike RA), crepitus, reduced range of motion, and joint enlargement (osteophytes, bony hypertrophy). Weight-bearing joints (knees, hips) and the hands (distal interphalangeal joints---Heberden nodes, proximal interphalangeal joints---Bouchard nodes, first carpometacarpal joints) are most commonly affected. Diagnosis is primarily clinical; X-rays show joint space narrowing, subchondral sclerosis, osteophytes, and subchondral cysts. Management includes patient education, weight loss, exercise, physical therapy, analgesics (acetaminophen, NSAIDs), intra-articular corticosteroid or hyaluronic acid injections, and joint replacement surgery (total knee or hip arthroplasty) for advanced disease. Prognosis is variable; joint replacement provides excellent outcomes for advanced disease.

%-------------------------------------------------------------
\section{Inflammatory Arthropathies}
%-------------------------------------------------------------

%-------------------------------------------------------------

%-------------------------------------------------------------

%-------------------------------------------------------------

%-------------------------------------------------------------

%-------------------------------------------------------------

%-------------------------------------------------------------

%-------------------------------------------------------------

\label{sec:Inflammatory Arthropathies}
%-------------------------------------------------------------%-------------------------------------------------------------

\subsection{Ankylosing Spondylitis}
\label{sec:Ankylosing Spondylitis}
Ankylosing spondylitis (AS) is a chronic inflammatory disease primarily affecting the axial skeleton (sacroiliac joints and spine), representing the prototypical spondyloarthropathy. It predominantly affects young men (20--40 years) and is strongly associated with HLA-B27 (present in 90\% of patients). The condition results from immune-mediated inflammation of the entheses and axial joints. The hallmark is sacroiliitis presenting with inflammatory back pain (insidious onset, morning stiffness >30 minutes, improvement with activity, onset before age 40). Clinical features include progressive spinal involvement leading to syndesmophyte formation, vertebral fusion ("bamboo spine"), and spinal rigidity; peripheral arthritis (hips, shoulders, knees), enthesitis (Achilles tendinitis, plantar fasciitis), anterior uveitis, and aortic regurgitation are associated features. Diagnosis involves clinical assessment, HLA-B27 testing, and imaging (MRI for early sacroiliitis, X-rays showing sacroiliac joint erosions and fusion, syndesmophytes). Management includes NSAIDs (first-line, used continuously), TNF inhibitors (adalimumab, etanercept, infliximab, certolizumab), and IL-17 inhibitors (secukinumab, ixekizumab) for inadequate response; physical therapy is essential. Prognosis is variable; early treatment improves outcomes.

\subsection{Gout}
\label{sec:Gout}
Gout is the most common inflammatory arthropathy, caused by deposition of monosodium urate crystals in joints and soft tissues due to chronic hyperuricemia (serum urate >6.8 mg/dL, the saturation point). It predominantly affects middle-aged men and postmenopausal women. The condition results from chronic hyperuricemia with risk factors including high-purine diet, alcohol (especially beer), obesity, chronic kidney disease, and certain medications (thiazide diuretics). Acute gouty arthritis classically presents as sudden, severe monoarthritis of the first metatarsophalangeal joint (podagra), with intense pain, erythema, warmth, and swelling; recurrent attacks may involve multiple joints. Chronic tophaceous gout causes tophi (urate deposits) in joints, cartilage, and soft tissues. Diagnosis is based on synovial fluid analysis showing needle-shaped, negatively birefringent crystals under polarized light microscopy. Management of acute attacks includes NSAIDs, colchicine (most effective within 12 hours of onset), or corticosteroids; long-term urate-lowering therapy (ULT) with allopurinol or febuxostat (xanthine oxidase inhibitors) or probenecid (uricosuric) is indicated for recurrent gout, tophi, or joint damage, targeting serum urate <6 mg/dL. Prognosis is excellent with effective urate-lowering therapy.

\subsection{Psoriatic Arthritis}
\label{sec:Psoriatic Arthritis}
Psoriatic arthritis (PsA) is discussed in Chapter~\ref{ch:skin} (Skin Diseases). It is an inflammatory arthropathy associated with psoriasis, classified as a seronegative spondyloarthropathy. Five clinical patterns include oligoarthritis, polyarthritis, arthritis mutilans, spondylitis, and dactylitis. Management includes NSAIDs, DMARDs (methotrexate), and biologics (TNF inhibitors, IL-17 inhibitors, IL-23 inhibitors, JAK inhibitors).

\subsection{Reactive Arthritis}
\label{sec:Reactive Arthritis}
Reactive arthritis is an aseptic inflammatory arthritis occurring 1--4 weeks after a genitourinary (usually \textit{Chlamydia trachomatis}) or gastrointestinal (usually \textit{Salmonella}, \textit{Shigella}, \textit{Yersinia}, \textit{Campylobacter}) infection. It is strongly associated with HLA-B27 and classified as a seronegative spondyloarthropathy. The condition results from an immune response to the initial infection triggering joint inflammation. The classic triad is arthritis, urethritis, and conjunctivitis ("can't see, can't pee, can't climb a tree"). Arthritis is typically asymmetric and oligoarticular, affecting large joints of the lower extremities; enthesitis, dactylitis, and oral ulcers are common. Most cases resolve within 3--12 months. Diagnosis is clinical with history of preceding infection. Management includes NSAIDs, intra-articular corticosteroids, and for chronic or severe disease, DMARDs (sulfasalazine) or TNF inhibitors; antibiotics are not helpful once arthritis has developed. Prognosis is generally good; most cases resolve within 3--12 months.

\subsection{Rheumatoid Arthritis}
\label{sec:Rheumatoid Arthritis}
Rheumatoid arthritis (RA) is a chronic, systemic autoimmune disease characterized by symmetric inflammatory synovitis, predominantly affecting the small joints of the hands (metacarpophalangeal and proximal interphalangeal joints) and feet. It affects approximately 0.5--1\% of the population, with a 3:1 female predominance. The condition results from autoantibody production (rheumatoid factor, anti-citrullinated peptide antibodies [anti-CCP]), immune complex deposition, and synovial inflammation leading to pannus formation and progressive joint destruction. Extra-articular manifestations include rheumatoid nodules, interstitial lung disease, vasculitis, scleritis, pericarditis, and accelerated atherosclerosis. Diagnosis requires clinical assessment, serological testing (RF, anti-CCP), and imaging (X-rays showing periarticular osteopenia, joint space narrowing, and erosions; ultrasound or MRI for early detection of synovitis). The 2010 ACR/EULAR classification criteria aid early diagnosis. Management follows a treat-to-target strategy with early initiation of DMARDs: methotrexate (first-line anchor drug), leflunomide, sulfasalazine, and hydroxychloroquine; biologics (TNF inhibitors, IL-6 inhibitors [tocilizumab, sarilumab], T-cell co-stimulation blocker [abatacept], anti-CD20 [rituximab]) and targeted synthetic DMARDs (JAK inhibitors [tofacitinib, baricitinib]) are used for inadequate response. Prognosis is variable but has improved significantly with early aggressive treatment.

%-------------------------------------------------------------
\section{Metabolic Bone Diseases}
%-------------------------------------------------------------

%-------------------------------------------------------------

%-------------------------------------------------------------

%-------------------------------------------------------------

%-------------------------------------------------------------

%-------------------------------------------------------------

%-------------------------------------------------------------

%-------------------------------------------------------------

\label{sec:Metabolic Bone Diseases}
%-------------------------------------------------------------%-------------------------------------------------------------

\subsection{Osteogenesis Imperfecta}
\label{sec:Osteogenesis Imperfecta}
Osteogenesis imperfecta (OI, "brittle bone disease") is a group of genetic disorders primarily caused by mutations in the COL1A1 or COL1A2 genes encoding type I collagen, leading to defective bone formation and increased bone fragility. Severity ranges from perinatal lethal (type II) to mild (type I, most common). The condition results from defective collagen production leading to brittle bones. Clinical features include recurrent fractures with minimal trauma, blue sclerae, dentinogenesis imperfecta (brittle, discolored teeth), hearing loss (conductive, sensorineural, or mixed), and joint hyperlaxity. Diagnosis is clinical, supported by genetic testing and bone densitometry. Management includes bisphosphonates (to reduce fracture frequency), physical therapy, bracing, surgical correction of deformities (intramedullary rodding), and dental management. Prognosis depends on the type; milder forms have a normal life expectancy.

\subsection{Osteoporosis}
\label{sec:Osteoporosis}
Osteoporosis is a systemic skeletal disease characterized by low bone mineral density (BMD) and deterioration of bone microarchitecture, resulting in increased bone fragility and fracture risk. It affects approximately 200 million people worldwide, predominantly postmenopausal women and the elderly. The condition results from an imbalance between bone resorption and bone formation, with risk factors including female sex, advancing age, low body weight, family history, smoking, excessive alcohol, low calcium and vitamin D intake, sedentary lifestyle, and medications (glucocorticoids, aromatase inhibitors, androgen deprivation therapy). Clinical features include fragility fractures commonly at the hip, vertebrae (compression fractures), and distal radius (Colles' fracture). Diagnosis is based on dual-energy X-ray absorptiometry (DXA): T-score $\leq$-2.5 at the hip or spine defines osteoporosis; T-score between -1 and -2.5 is osteopenia. The FRAX tool estimates 10-year fracture risk. Management includes calcium (1000--1200 mg/day) and vitamin D (800--1000 IU/day), weight-bearing exercise, bisphosphonates (alendronate, zoledronic acid, first-line), denosumab (RANKL inhibitor), teriparatide and abaloparatide (parathyroid hormone analogs), and romosozumab (anti-sclerostin antibody). Prognosis is improved with early detection and treatment; fracture risk is significantly reduced with appropriate therapy.

\subsection{Paget's Disease of Bone}
\label{sec:Paget's Disease of Bone}
Paget's disease of bone is a chronic disorder of bone remodeling characterized by excessive osteoclastic bone resorption followed by disordered, chaotic new bone formation, resulting in enlarged, structurally abnormal, weakened bones. It is the second most common metabolic bone disease after osteoporosis, affecting 1--3\% of adults over 55; the pelvis, spine, skull, and femur are most commonly involved. The condition results from abnormal osteoclast activity with increased bone turnover. Most patients are asymptomatic, discovered incidentally by elevated alkaline phosphatase or radiographic findings. Clinical features include bone pain, skeletal deformity, pathologic fractures, and neurological complications (compression of cranial nerves, spinal stenosis, hearing loss from skull base involvement). Diagnosis is based on elevated serum alkaline phosphatase, characteristic radiographic changes (osteolytic, mixed, or osteosclerotic phases, cotton-wool appearance of the skull), and bone scan showing increased uptake. Management with bisphosphonates (zoledronic acid most effective) inhibits osteoclastic activity and normalizes bone turnover. Prognosis is generally good with treatment.

\subsection{Rickets and Osteomalacia}
\label{sec:Rickets and Osteomalacia}
Rickets (in children) and osteomalacia (in adults) are disorders of impaired bone mineralization due to vitamin D deficiency or abnormalities in vitamin D metabolism. The condition results from inadequate vitamin D (produced in the skin upon UV exposure and obtained from diet), which is essential for intestinal calcium absorption, leading to soft, weak bones. Clinical features in children (rickets) include growth retardation, skeletal deformities (bowed legs, knock knees, rachitic rosary, Harrison's groove, frontal bossing), dental abnormalities, and hypocalcemic seizures. Clinical features in adults (osteomalacia) include diffuse bone pain, muscle weakness, increased fracture risk, and pathologic fractures. Diagnosis is based on serum 25-hydroxyvitamin D level (<20 ng/mL indicates deficiency), elevated alkaline phosphatase, low calcium and phosphate, and radiographic findings (widened growth plates, Looser's zones in adults). Management is vitamin D supplementation (ergocalciferol or cholecalciferol) and calcium. Prognosis is excellent with adequate vitamin D replacement.

%-------------------------------------------------------------
\section{Other Bone and Joint Conditions}
%-------------------------------------------------------------

%-------------------------------------------------------------

%-------------------------------------------------------------

%-------------------------------------------------------------

%-------------------------------------------------------------

%-------------------------------------------------------------

%-------------------------------------------------------------

%-------------------------------------------------------------

\label{sec:Other Bone and Joint Conditions}
%-------------------------------------------------------------%-------------------------------------------------------------

\subsection{Achondroplasia}
\label{sec:Achondroplasia}
Achondroplasia is the most common form of disproportionate short stature, caused by a gain-of-function mutation in FGFR3 (most commonly G380R), with an incidence of 1 in 25,000 live births. The condition results from constitutive activation of FGFR3, which negatively regulates endochondral ossification in the growth plate, leading to impaired long bone growth. Clinical features include rhizomelic limb shortening (proximal segments more affected), macrocephaly with frontal bossing, midface hypoplasia, trident hand, genu varum, thoracolumbar kyphosis, and foramen magnum stenosis. Intelligence is normal. Complications include cervicomedullary compression, hydrocephalus, spinal stenosis, and obesity. Diagnosis is based on clinical and radiographic findings with FGFR3 genetic confirmation. Management involves monitoring for hydrocephalus, decompression surgery for spinal stenosis, limb lengthening procedures (controversial), and vosoritide (CNP analog, FDA-approved in 2021) to improve growth velocity. Prognosis is generally good; life expectancy is reduced by 10 years primarily due to cardiovascular complications and spinal stenosis.

\subsection{Fibrous Dysplasia}
\label{sec:Fibrous Dysplasia}
Fibrous dysplasia is a benign bone disorder where normal bone is replaced by fibrous tissue and immature woven bone, caused by postzygotic activating mutations in the GNAS gene. It may be monostotic (single bone, 70--80\%) or polyostotic (multiple bones); the ribs, femur, tibia, and skull are commonly affected. The condition results from GNAS mutation leading to abnormal bone formation. McCune-Albright syndrome is a variant combining polyostotic fibrous dysplasia with caf\'{e}-au-lait spots and endocrine hyperfunction (precocious puberty). Clinical features include bone pain, deformity, or pathologic fracture. Diagnosis is by X-ray showing ground-glass appearance and thinning of the cortex. Management is symptomatic: observation, bisphosphonates for pain, and surgical correction of deformities or fractures. Prognosis is generally good; malignant transformation is rare.

\subsection{Fractures}
\label{sec:Fractures}
A fracture is a break in the continuity of bone. Fractures are classified as open (compound, with skin disruption) or closed (simple), and by pattern (transverse, oblique, spiral, comminuted, greenstick, impacted, avulsion). Common fracture types include distal radius (Colles' fracture), hip (intertrochanteric or subcapital femoral neck), vertebral compression, scaphoid, and clavicle fractures. The condition results from mechanical force exceeding bone strength. Clinical features include pain, deformity, swelling, and loss of function. Diagnosis is by X-rays (at least two views); CT for complex fractures, MRI for occult fractures and ligamentous injury. Management aims for reduction (realigning fracture fragments), fixation (maintaining alignment during healing), and rehabilitation; closed reduction with casting is appropriate for stable, minimally displaced fractures; open reduction with internal fixation (ORIF) is indicated for displaced, unstable, or intra-articular fractures. Prognosis is generally good; hip fractures in the elderly are orthopedic emergencies requiring surgical fixation within 24--48 hours to reduce mortality and complications.

\subsection{Kyphosis}
\label{sec:Kyphosis}
Kyphosis is excessive anterior curvature of the thoracic spine (normal kyphosis 20--45$^\circ$). Causes include Scheuermann's disease (adolescent, with vertebral wedging), osteoporotic compression fractures (elderly women---"dowager's hump"), ankylosing spondylitis, congenital anomalies, and postural causes. The condition results from the specific underlying cause leading to vertebral wedging or collapse. Scheuermann's kyphosis presents in adolescence with rigid thoracic kyphosis and is diagnosed by the presence of >5$^\circ$ of wedging in three or more consecutive vertebrae on X-ray. Clinical features include visible curvature, back pain, and stiffness. Diagnosis is based on physical examination and X-ray. Management ranges from observation and physical therapy for mild cases to bracing and surgery for severe deformities. Prognosis is generally good; most patients function well with conservative management.

\subsection{Scoliosis}
\label{sec:Scoliosis}
Scoliosis is abnormal lateral curvature of the spine (Cobb angle >10$^\circ$). Adolescent idiopathic scoliosis (AIS) is the most common form (80\% of cases), typically presenting in girls during the adolescent growth spurt, with a right thoracic curve. Neuromuscular scoliosis (from cerebral palsy, muscular dystrophy), congenital scoliosis (vertebral anomalies), and syndromic scoliosis (Marfan, Down syndrome) are other forms. The condition results from abnormal spinal development or neuromuscular imbalance. Clinical features include asymmetrical trunk, shoulder height discrepancy, and rib hump on forward bending. Diagnosis is based on physical examination and X-ray with Cobb angle measurement. Management depends on severity: mild curves are monitored; moderate curves (25--40$^\circ$) are treated with bracing; severe curves (>40--50$^\circ$) require surgical correction with posterior spinal instrumentation and fusion. Prognosis is generally favorable with appropriate management; complications include chronic back pain and cardiopulmonary compromise in severe curves.

%-------------------------------------------------------------
\section{Soft Tissue and Overuse Conditions}
%-------------------------------------------------------------

%-------------------------------------------------------------

%-------------------------------------------------------------

%-------------------------------------------------------------

%-------------------------------------------------------------

%-------------------------------------------------------------

%-------------------------------------------------------------

%-------------------------------------------------------------

\label{sec:Soft Tissue and Overuse Conditions}
%-------------------------------------------------------------%-------------------------------------------------------------

\subsection{Adhesive Capsulitis}
\label{sec:Adhesive Capsulitis}
Adhesive capsulitis (frozen shoulder) is a chronic condition characterized by progressive pain and global restriction of passive and active range of motion of the glenohumeral joint, affecting 2--5\% of the population. The condition results from synovial inflammation followed by capsular thickening, pericapsular contracture, and dense fibrotic adhesions leading to severe loss of joint volume. Risk factors include diabetes mellitus (most significant), thyroid dysfunction, cerebrovascular accidents, prolonged shoulder immobilization, and autoimmune disease. Clinical progression follows three overlapping phases: freezing (painful, 2--9 months), frozen (stiffness spredominates, 4--12 months), and thawing (gradual motion recovery, 5--24 months). Diagnosis is clinical based on equal restriction of active and passive external rotation; radiographs rule out glenohumeral osteoarthritis. Management involves physical therapy, intra-articular corticosteroid injections, oral NSAIDs, and manipulation under anesthesia or capsular release for recalcitrant cases.

\subsection{Bursitis}
\label{sec:Bursitis}
Bursitis is inflammation of a bursa, a synovial fluid--filled sac that reduces friction between bones and overlying soft tissues (tendons, muscles, skin). The most commonly affected bursae include the subacromial (shoulder), olecranon (elbow), prepatellar (knee---"housemaid's knee"), and trochanteric (hip---"weaver's bottom"). Risk factors include repetitive motion or prolonged pressure, overuse, trauma, inflammatory conditions (gout, rheumatoid arthritis), and infection (septic bursitis). The condition results from mechanical irritation or infection of the bursa. Clinical features include localized pain, swelling, tenderness, and sometimes erythema and warmth over the affected bursa; the overlying joint typically has a full range of motion, distinguishing bursitis from arthritis. Septic bursitis is more common in the olecranon and prepatellar bursae, often caused by \textit{S.~aureus}, presenting with more intense inflammation, fever, and systemic symptoms. Diagnosis is primarily clinical; ultrasound can confirm bursal fluid collection. Aspiration of bursal fluid is indicated when septic bursitis is suspected. Management of aseptic bursitis includes rest, ice, NSAIDs, and corticosteroid injection; septic bursitis requires bursal aspiration, incision and drainage if needed, and antibiotics. Prognosis is excellent with appropriate treatment.

\subsection{Calcium Pyrophosphate Deposition Disease}
\label{sec:Calcium Pyrophosphate Deposition Disease}
Calcium pyrophosphate deposition disease (CPPD, formerly pseudogout) is a crystal-induced arthropathy resulting from the deposition of calcium pyrophosphate dihydrate (CPP) crystals in articular cartilage (chondrocalcinosis) and synovium. The condition occurs predominantly in elderly individuals ($>60$ years) and is associated with hyperparathyroidism, hemochromatosis, hypomagnesemia, and hypothyroidism. Clinical features include acute monoarthritis or oligoarthritis resembling gout, most commonly affecting the knee, wrist, or symphysis pubis, with swelling, warmth, severe pain, and erythema. Diagnosis is confirmed by synovial fluid aspiration showing weakly positively birefringent, rhomboid-shaped CPP crystals under polarized light microscopy; radiographs demonstrate chondrocalcinosis (linear calcification of articular cartilage). Management of acute attacks includes intra-articular corticosteroid injection, oral NSAIDs, or colchicine. Prognosis is good for acute control.

\subsection{Dupuytren Contracture}
\label{sec:Dupuytren Contracture}
Dupuytren contracture is a progressive fibroproliferative disorder of the palmar fascia resulting in nodule formation and fixed flexion deformities of the digits. The condition results from abnormal myofibroblast proliferation and excessive collagen type III deposition, leading to thickening and shortening of the palmar aponeurosis. Risk factors include Northern European ancestry, male sex, advanced age, positive family history, diabetes mellitus, heavy alcohol consumption, and smoking. Clinical features begin with painless palmar skin pitting and subcutaneous nodules (most commonly over the 4th and 5th rays), progressing to dense fascial cords and fixed flexion contractures at the metacarpophalangeal (MCP) and proximal interphalangeal (PIP) joints (positive Tabletop test). Diagnosis is clinical. Management includes observation for mild non-functional cases, collagenase Clostridium histolyticum injections, needle aponeurotomy, or subtotal palmar fasciectomy for functional impairment ($>30^\circ$ MCP contracture).

\subsection{Hallux Valgus}
\label{sec:Hallux Valgus}
Hallux valgus is a structural deformity of the first ray of the foot characterized by lateral deviation of the hallux (great toe) and medial displacement of the first metatarsal head. The condition results from biomechanical instability of the first metatarsophalangeal (MTP) joint, often exacerbated by constrictive footwear, ligamentous laxity, flat feet (pes planus), or genetic predisposition. Clinical features include a painful exostosis (bunion) over the medial eminence of the first metatarsal head, lateral crowding of second and third toes, metatarsalgia, and difficulty fitting shoes. Diagnosis is confirmed by weight-bearing foot radiographs showing a first MTP angle $>15^\circ$ and an intermetatarsal angle $>9^\circ$. Management begins with wide-toebox shoes, orthotics, and night splints; surgical osteotomy and soft tissue balancing are indicated for severe refractory pain.

\subsection{Morton Neuroma}
\label{sec:Morton Neuroma}
Morton neuroma is a benign perineural fibrotic enlargement of the interdigital plantar nerve, most commonly affecting the third intermetatarsal space between the third and fourth metatarsal heads. The condition results from chronic mechanical compression and repetitive irritation of the nerve beneath the deep transverse metatarsal ligament, leading to epineural and endoneural sclerosis. Clinical features include sharp, burning forefoot pain radiating to the toes, exacerbated by weight-bearing and narrow footwear, often described as 'walking on a marble.' Physical examination reveals a positive Mulder sign (palpable click with forefoot squeeze) and localized intermetatarsal tenderness. Diagnosis is clinical, supported by foot ultrasound or MRI if equivocal. Management includes wide footwear, metatarsal pads, corticosteroid injections, and surgical neurectomy for recalcitrant symptoms.

\subsection{Plantar Fasciitis}
\label{sec:Plantar Fasciitis}
Plantar fasciitis is the most common cause of heel pain in adults, resulting from overuse injury and degeneration of the plantar fascia, a thick band of connective tissue originating from the medial tubercle of the calcaneus and inserting into the toes. It affects approximately 10\% of the population at some point in their life, with peak incidence between ages 40 and 60. The condition results from repetitive microtrauma and degeneration of the plantar fascia. Risk factors include obesity, flat feet, high arches, prolonged standing, running, and tight Achilles tendon. The classic presentation is sharp heel pain with the first steps in the morning or after rest ("start-up pain"), which improves with activity but may worsen with prolonged standing; pain is maximal at the medial heel. Diagnosis is clinical; imaging (X-ray, ultrasound, or MRI) is rarely needed but may show calcaneal spur (present in 50\% of cases but often asymptomatic) or plantar fascia thickening. Management includes rest, activity modification, stretching exercises (calf and plantar fascia), ice massage, orthotics, and night splints; corticosteroid injection may provide short-term relief but should be limited; extracorporeal shockwave therapy is an alternative for refractory cases. Prognosis is excellent; most cases resolve within 6--18 months with conservative management.

\subsection{Trigger Finger}
\label{sec:Trigger Finger}
Trigger finger (stenosing tenosynovitis) is a common tenosynovial disorder characterized by painful catching, snapping, or locking of the digit during flexion and extension. The condition results from a mismatch between the diameter of the flexor digitorum tendon and the narrowed first annular (A1) pulley at the metacarpal head, driven by repetitive microtrauma, tendon nodule formation, and fibrocartilaginous metaplasia. Risk factors include diabetes mellitus, rheumatoid arthritis, thyroid dysfunction, and repetitive hand gripping. Clinical features include localized pain and a palpable tender nodule over the palmar MTP joint, with mechanical clicking or locking of the digit that may require passive manual extension to unlock. Diagnosis is clinical. Management includes activity modification, night splinting, local corticosteroid injection into the tendon sheath (highly effective), and surgical A1 pulley release for refractory cases.

